% !TEX program = xelatex
% 공통수학2 · Ⅰ. 도형의 방정식 · 3. 도형의 이동 · 02 대칭이동 · 1/2차시
\documentclass[11pt,a4paper]{article}
\usepackage{kotex}
\usepackage{iftex}
\ifPDFTeX\else
  \setmainhangulfont{Noto Serif CJK KR}
  \setsanshangulfont{Noto Sans CJK KR}
\fi
\usepackage[margin=2.1cm]{geometry}
\usepackage{parskip}
\usepackage{enumitem}
\usepackage{array}
\usepackage{tabularx}
\usepackage{booktabs}
\usepackage{xcolor}
\usepackage{amssymb}
\usepackage{tikz}
\usepackage[most]{tcolorbox}
\usepackage{fancyhdr}
\usepackage{titlesec}

\definecolor{goldc}{HTML}{B07C22}
\definecolor{goldstrong}{HTML}{8A5F16}
\definecolor{indigoc}{HTML}{2B3B57}
\definecolor{indigosoft}{HTML}{6E82A6}
\definecolor{linec}{HTML}{D6DACB}
\definecolor{surface2}{HTML}{E4E8DC}
\definecolor{notebg}{HTML}{FBF3E3}
\definecolor{inksoft}{HTML}{52604F}

\titleformat{\section}{\normalfont\sffamily\bfseries\large\color{indigoc}}{\thesection}{0.6em}{}
\titlespacing{\section}{0pt}{16pt}{8pt}
\newtcolorbox{ideabox}{enhanced, breakable, colback=white, colframe=goldc, boxrule=0.5pt, leftrule=3pt, arc=2pt, left=10pt, right=10pt, top=8pt, bottom=8pt}
\newtcolorbox{calloutbox}{enhanced, breakable, colback=white, colframe=indigosoft, boxrule=0.5pt, leftrule=3pt, arc=2pt, left=10pt, right=10pt, top=7pt, bottom=7pt}
\newtcolorbox{notebox}{enhanced, breakable, colback=notebg, colframe=goldc, boxrule=0.5pt, leftrule=3pt, arc=2pt, left=10pt, right=10pt, top=7pt, bottom=7pt}
\newtcolorbox{answerbox}{enhanced, breakable, colback=white, colframe=linec, boxrule=0.5pt, arc=2pt, left=10pt, right=10pt, top=7pt, bottom=7pt}
\newcommand{\lessonbanner}[3]{%
  \begin{tcolorbox}[enhanced, colback=indigoc, colframe=indigoc, boxrule=0pt, arc=0pt,
    left=14pt, right=14pt, top=14pt, bottom=10pt, borderline south={2.4pt}{0pt}{goldc}, fontupper=\color{white}]
    {\sffamily\footnotesize\color{white!85!black} #1}\\[4pt]{\sffamily\bfseries\Large #2}\\[3pt]{\sffamily\small\color{white!88!black} #3}
  \end{tcolorbox}}
\pagestyle{fancy}\fancyhf{}\renewcommand{\headrulewidth}{0pt}
\fancyfoot[L]{\footnotesize\color{inksoft} 공통수학2 · 2026학년도 2학기 · 지평선고등학교}
\fancyfoot[R]{\footnotesize\color{inksoft} 교사용 지도안 -- 배포 금지}

\begin{document}
\lessonbanner{공통수학2 · Ⅰ. 도형의 방정식 · 3. 도형의 이동}{02 대칭이동 --- 1/2차시}{점의 대칭이동 · 교사용 지도안 (2026학년도 2학기)}
\vspace{10pt}

\noindent\begin{tabularx}{\linewidth}{@{}>{\bfseries\color{indigoc}}p{2.6cm}X@{}}
\toprule
대상 & 고1 · 2개 반 · 반당 13명 \\
차시 & 50분 (도입 10 · 전개 30 · 정리 10) \\
성취기준 & 원점, $x$축, $y$축, 직선 $y=x$에 대한 대칭이동을 탐구하고, 실생활과 연결하여 문제를 해결할 수 있다. \\
선행 학습 & 14차시 --- 평행이동 \\
\bottomrule
\end{tabularx}

\section{학습 목표 \& Big Idea}
\begin{ideabox}
\textbf{\color{goldstrong}영속적 이해 (Big Idea)}\\
대칭이동은 `거울에 비추는' 것이다. 어느 거울(축)에 비추느냐에 따라 좌표의 어느 부호가 뒤집히는지가 정확히 결정된다.
\vspace{4pt}
\small\color{inksoft} 핵심개념: \textbf{대칭} \quad\textbullet\quad 핵심개념: \textbf{부호의 반전} \quad\textbullet\quad 일반화: \textbf{거울(기준)이 달라지면 뒤집히는 좌표도 달라진다}
\end{ideabox}
\begin{calloutbox}
\textbf{차시 학습 목표}
\begin{itemize}[leftmargin=1.4em, itemsep=2pt, topsep=4pt]
  \item 점을 $x$축, $y$축, 원점, 직선 $y=x$에 대하여 대칭이동한 점의 좌표를 구할 수 있다.
\end{itemize}
\end{calloutbox}

\section{질문의 연결}
\begin{tabularx}{\linewidth}{@{}>{\bfseries\color{indigoc}\small}p{2.4cm}X@{}}
사실적 질문 & 점 $(x,y)$를 $x$축, $y$축, 원점, 직선 $y=x$에 대하여 대칭이동하면 좌표는 각각 어떻게 바뀔까? \\[6pt]
개념적 질문 & 왜 직선 $y=x$에 대한 대칭이동만 좌표가 `뒤바뀌고'(순서 교환), 나머지는 부호만 바뀔까? \\[6pt]
논쟁적 질문 & 마천루처럼 좌우 대칭인 건축은 왜 아름답고 안정적으로 느껴질까? 대칭이 항상 최선의 디자인일까? \\
\end{tabularx}

\section{수업 흐름}
\begin{tabularx}{\linewidth}{@{}>{\centering\arraybackslash}p{1.6cm}X>{\raggedright\arraybackslash\small}p{3.1cm}@{}}
\toprule
\textbf{단계} & \textbf{활동 \& 발문} & \textbf{ATL / VTR} \\
\midrule
\textcolor{goldstrong}{\textbf{도입}}\newline\footnotesize 10분 &
\textbf{마음공부 (2분)} --- 유무념 대조.\newline
\textbf{생각 열기 (8분)} --- 만자창 무늬를 좌표평면에 나타내고, 점 A$(1,2)$를 $x$축·$y$축·원점에 대하여 대칭이동한 점 B, C, D를 직접 그려서 좌표를 말해 보게 한다. &
ATL: 사고(패턴 관찰)\newline VTR: See-Think-Wonder \\
\addlinespace
\textcolor{goldstrong}{\textbf{전개}}\newline\footnotesize 30분 &
\textbf{개념 형성 (15분)} --- $x$축, $y$축, 원점에 대한 대칭이동 공식을 `PP$'$의 중점' 논리로 유도 $\to$ 직선 $y=x$에 대한 대칭이동은 (①중점이 $y=x$ 위에 있다, ②PP$'\perp y=x$) 두 조건의 연립으로 유도 $\to$ 문제 1.\newline
\textbf{심화 (15분)} --- 직선 $y=-x$에 대한 대칭이동, 점 $(a,b)$에 대한 대칭이동(대칭의 중심 = 중점)을 같은 논리로 확장 --- 학생이 스스로 유도해 보게 한다. &
ATT: 촉진자 역할\newline ATL: 사고·일반화 \\
\addlinespace
\textcolor{goldstrong}{\textbf{정리}}\newline\footnotesize 10분 &
\textbf{개념 연결 질문 (5분)} --- ``네 가지 대칭이동 중 부호가 바뀌는 것과 좌표가 뒤바뀌는 것, 어떻게 구별할까?''\newline
\textbf{성찰 (5분)} --- 감사 일기 1문장 + 다음 차시 예고(도형의 대칭이동). &
마음공부 · 감사 일기 \\
\bottomrule
\end{tabularx}

\begin{calloutbox}
\textbf{지도상의 유의점} --- 원점에 대한 대칭이동과 직선에 대한 대칭이동을 혼동하지 않도록 주의한다. 네 가지 공식을 암기시키기보다 `PP$'$의 중점이 대칭축(또는 대칭점) 위에 있다'는 하나의 원리에서 모두 유도됨을 강조한다.
\end{calloutbox}

\section{형성평가 \& 답안}
\begin{answerbox}
\textbf{문제 1} \quad 점 $(3,-5)$를 대칭이동\\
\small ⑴ $x$축 \quad ⑵ $y$축 \quad ⑶ 원점 \quad ⑷ 직선 $y=x$\\[4pt]
\textbf{\color{goldstrong}답} \; ⑴ $(3,5)$ \quad ⑵ $(-3,-5)$ \quad ⑶ $(-3,5)$ \quad ⑷ $(-5,3)$
\end{answerbox}

\section{성취수준 (이 차시 범위)}
\begin{tabularx}{\linewidth}{@{}>{\bfseries\color{goldstrong}}p{0.9cm}X@{}}
\toprule
A & 원점, $x$축, $y$축, 직선 $y=x$에 대한 대칭이동을 탐구하고 설명할 수 있다. \\
B & 대칭이동을 이해하고 대칭이동한 점의 좌표를 구할 수 있다. \\
C & 대칭이동한 점의 좌표를 구할 수 있다. \\
D & 대칭이동한 점의 좌표를 구할 수 있다(안내됨). \\
E & $x$축, $y$축에 대하여 대칭이동한 점의 좌표를 구할 수 있다. \\
\bottomrule
\end{tabularx}

\section{다음 차시}
\begin{calloutbox}
\textbf{02 대칭이동 --- 2/2차시.} 점의 대칭이동 공식을 그대로 확장하여, 방정식이 나타내는 도형 전체를 대칭이동하는 방법을 다룬다.
\end{calloutbox}
\end{document}
